\chapter{Newton's method}
\label{ch:newton}
\impl{src/nonlinear\_solvers.jl}

\section{The iteration}

Newton--Raphson linearizes $\Rvec(\Uvec) = \vect{0}$ about the current iterate:
\begin{equation}
  \Kmat(\Uvec_k)\, \Delta\Uvec_k = \Rvec(\Uvec_k),
  \qquad
  \Uvec_{k+1} = \Uvec_k + \alpha_k \Delta\Uvec_k ,
  \label{eq:newton}
\end{equation}
with $\Kmat$ the tangent \eqref{eq:tangent} for quasi-static problems and
$\Keff$ of \eqref{eq:keff} for implicit dynamics.  With an exact tangent, an
exact linear solve, and $\alpha_k = 1$, convergence near the solution is
quadratic.  Each of those three assumptions is relaxed somewhere in Carina, and
each relaxation is the subject of a later chapter: approximate tangents in
Section~\ref{sec:tangent-strategies}, inexact linear solves in
Chapter~\ref{ch:inexact}, and $\alpha_k < 1$ immediately below.

\section{Globalization by line search}
\label{sec:linesearch}

Far from the solution the full Newton step can be worse than no step at all.
At large $\Delta t$ or a large load increment the predictor may sit far enough
from equilibrium that $\Uvec_k + \Delta\Uvec_k$ inverts elements, producing
$J \le 0$ and a non-finite residual.

Carina globalizes with Armijo backtracking on the merit function
\begin{equation}
  m(\alpha) = \tfrac12 \norm{\Rvec(\Uvec_k + \alpha \Delta\Uvec_k)}^2 ,
\end{equation}
accepting the first $\alpha \in \{1, \tau, \tau^2, \dots\}$ satisfying the
sufficient-decrease condition
\begin{equation}
  m(\alpha) \le \left(1 - 2 c\, \alpha\right) m(0),
  \label{eq:armijo}
\end{equation}
with backtrack factor $\tau = \tfrac12$ and $c = 10^{-4}$ by default.
Equation~\eqref{eq:armijo} is the standard Armijo condition specialized to this
merit function: since $\Delta\Uvec_k$ solves $\Kmat\Delta\Uvec = \Rvec$, the
directional derivative is $m'(0) = -\norm{\Rvec}^2 = -2m(0)$, and substituting
into $m(\alpha) \le m(0) + c\,\alpha\, m'(0)$ gives \eqref{eq:armijo} directly.

\begin{algorithm}[h]
\caption{Newton step with Armijo backtracking}
\label{alg:newton}
\begin{algorithmic}[1]
  \Require iterate $\Uvec_k$, residual $\Rvec_k$, tangent $\Kmat_k$,
           parameters $c$, $\tau$, budget $j_{\max}$
  \State $\eta_k \gets$ \Call{ForcingTerm}{$\cdot$}
    \Comment{Algorithm~\ref{alg:forcing}; $\eta_k = \varepsilon_{\mathrm{lin}}$ if disabled}
  \State solve $\Kmat_k \Delta\Uvec_k = \Rvec_k$ to relative tolerance $\eta_k$
  \State $m_0 \gets \tfrac12\norm{\Rvec_k}^2$,\quad $\alpha \gets 1$
  \For{$j = 1$ \textbf{to} $j_{\max}$}
    \State evaluate $\Rvec$ at $\Uvec_k + \alpha\Delta\Uvec_k$;\quad
           $m \gets \tfrac12\norm{\Rvec}^2$
    \If{$m \le (1 - 2c\alpha)\,m_0$}
      \State \Return $\Uvec_k + \alpha\Delta\Uvec_k$
        \Comment{Armijo satisfied}
    \EndIf
    \State $\alpha \gets \tau\alpha$
  \EndFor
  \State restore the pre-step material state
    \Comment{trials advanced it; see Remark~\ref{rem:restore}}
  \State \Return $\Uvec_k + \Delta\Uvec_k$
    \Comment{budget exhausted; outer failure machinery takes over}
\end{algorithmic}
\end{algorithm}

The line search is \emph{on by default}.  It costs roughly one extra residual
evaluation per iteration when the full step is already good, since $\alpha = 1$
is then accepted immediately, and it is the only thing standing between a bad
predictor and an inverted element.

\begin{remark}
\label{rem:restore}
If the backtracking budget is exhausted without satisfying
\eqref{eq:armijo}, Carina restores the pre-step material state and accepts the
full step, leaving the outer failure machinery of Section~\ref{sec:adaptive} to
handle it.  Restoring the state matters: the trial evaluations have already
advanced the internal variables of a path-dependent model, and accepting a step
on top of a state advanced by a rejected trial would silently corrupt the
history.
\end{remark}

\begin{remark}
An inexact Newton direction (Chapter~\ref{ch:inexact}) remains a descent
direction for $m$ as long as the linear residual is smaller than $\norm{\Rvec}$
--- that is, whenever the forcing term satisfies $\eta_k < 1$, which the clamp
of Section~\ref{sec:ew-safeguards} enforces.  The line search therefore
composes with inexact solves without modification.
\end{remark}

\section{Status tests}
\label{sec:status}

Termination is a \emph{tree} of composable tests rather than a fixed criterion,
following the NOX pattern.  Leaves test the residual ($\norm{\Rvec} <
\varepsilon_a$ or $\norm{\Rvec}/\norm{\Rvec_0} < \varepsilon_r$), the update
($\norm{\Delta\Uvec}$, absolute or relative to $\norm{\Uvec}$), the iteration
count, finiteness, divergence, or stagnation; interior nodes combine them with
\textsc{And}\,/\,\textsc{Or}.

Two properties of the combination rules are worth stating because algorithms
elsewhere depend on them.  In an \textsc{Or} node, \textsc{Converged} beats
\textsc{Failed} when both occur on the same iteration.  And a
\emph{finite-value} test is appended to every tree automatically, so a
\textsc{NaN} residual always terminates the solve regardless of what the user
wrote.

The tree is not merely a convenience.  Section~\ref{sec:ew-safeguards} needs
the residual norm the solve is actually aiming at, and reads it off this tree
--- taking the loosest child of an \textsc{Or} and the tightest of an
\textsc{And} --- rather than from a nominal tolerance that a user-supplied tree
may have overridden.
